\documentclass{article}
\usepackage[margin=1in]{geometry}
\usepackage{longtable}
\begin{document}

\textbf{Legacy note:} This split-by-person draft complements the canonical roadmap in \texttt{Documentation/SRS.tex}.\\

\section*{Phase 7 Plan - Person 2 (Person B)}

\textbf{Source of truth:} \texttt{Documentation/SRS.tex}, Section 11.5, ``Phase 7: Advanced Research and Scale-Up''.\\
\textbf{Interpretation note:} Person 2 owns the advanced-research track after Phase 6. \textbf{Assume Person 1 already has the data, replay artifacts, and operational ML plumbing available.} Person 2 should spend time on research questions, new feature families, model experiments, ablations, and thesis-quality outputs rather than data engineering or collector operations.\\
\textbf{Implementation note:} Every advanced experiment must remain reproducible from saved datasets, hashes, and model versions. Research novelty does not excuse leakage or weak evaluation discipline.

\textbf{Phase duration:} after v1 stability\\
\textbf{Phase goal:} Expand the system into thesis-grade research depth on top of a stable data and runtime foundation.

\section*{Owned Deliverables (From SRS)}
\begin{enumerate}
\item Add graph-derived features and cluster persistence analysis.
\item Experiment with marked Hawkes or TCN sequence models.
\item Run the H1 observability / Goodhart study.
\item Produce ablation studies and thesis-quality result tables and figures.
\end{enumerate}

\section*{Locked Implementation Decisions}
\begin{itemize}
\item Person 2 should assume Person 1 provides stable datasets, restore guarantees, and artifact hashes.
\item Graph and temporal models must be compared against the Phase 6 baseline, not presented in isolation.
\item Observability / Goodhart work must be grounded in measurable operator-facing proxies, not vague discussion.
\item Thesis tables and figures must be reproducible from saved evaluation artifacts.
\item Expensive experiments must be prioritized by likely information gain, not curiosity alone.
\end{itemize}

\section*{Concrete Execution Breakdown}
\begin{longtable}{p{0.28\linewidth}p{0.32\linewidth}p{0.30\linewidth}}
\textbf{Work Item} & \textbf{Artifact} & \textbf{Done Condition} \\
Graph features & Graph / cluster feature set & Wallet-cluster persistence and relationship features are computed reproducibly from retained datasets \\
Advanced temporal models & Hawkes / TCN experiment track & At least one sequence-oriented family is evaluated against the Phase 6 baseline \\
Goodhart / observability study & H1 research memo & The system's observability incentives and possible gaming behaviors are analyzed with evidence \\
Ablations and thesis figures & Reproducible tables and figures & Results can be regenerated from saved datasets, model versions, and scripts \\
\end{longtable}

\section*{Execution Sequence (No Day Timeline)}
\textbf{Block 1: Data-ready research setup}
\begin{itemize}
\item Accept the stable replay datasets and hashes from Person 1.
\item Freeze one research index of windows, categories, and evaluation scopes.
\item Build the first graph-derived features from the retained data instead of asking Person 1 for new collection logic.
\end{itemize}

\textbf{Block 2: Advanced model experiments}
\begin{itemize}
\item Train graph-aware or cluster-aware variants on the stable Phase 6 baseline contract.
\item Run marked Hawkes and/or TCN experiments using the same label and split discipline as Phase 6.
\item Reject any apparent gain that disappears under strict holdouts or ablations.
\end{itemize}

\textbf{Block 3: Goodhart / observability study}
\begin{itemize}
\item Study whether alerting, scoring, or evidence policies create incentives that distort measured success.
\item Quantify observability blind spots, gaming risk, and metric brittleness.
\item Tie the findings back to concrete operator or thesis implications.
\end{itemize}

\textbf{Block 4: Thesis-quality reporting}
\begin{itemize}
\item Produce ablation tables, comparison figures, and methodology notes suitable for a thesis or paper appendix.
\item Package every key figure so it can be regenerated from versioned artifacts.
\item Hand Person 1 only the artifact set and deployment-relevant findings, not new infra burden.
\end{itemize}

\section*{Handoffs to Person 1}
\begin{itemize}
\item Versioned advanced-model artifacts and exact dataset hashes.
\item Comparison reports versus the Phase 6 baseline.
\item Any deployment-relevant warnings from the Goodhart / observability study.
\item Thesis figures and tables that rely on retained operational datasets.
\end{itemize}

\section*{Gate 7 Checklist (Person 2 Ownership)}
\begin{itemize}
\item Graph-derived or cluster-persistence features exist and are reproducible.
\item At least one advanced temporal model family is evaluated against the Phase 6 baseline.
\item The H1 observability / Goodhart study is written with evidence.
\item Thesis-quality ablation tables and figures can be regenerated from saved artifacts.
\end{itemize}

\end{document}
